\frfilename{mt2.tex} 
\versiondate{24.2.16} 
 
\noindent{\bf Pendahuluan Jilid 2} 
 
\medskip 
 
Untuk jilid kedua ini saya telah memilih tujuh topik untuk menelusuri
wawasan dan tantangan yang ditawarkan teori ukuran. Beberapa di antaranya,
seperti Teorema Radon-Nikod\'ym (Bab 23), mutlak diperlukan untuk memahami
struktur bidang ini; yang lain, seperti analisis Fourier (Bab 28) dan
pembahasan ruang fungsi (Bab 24), memperlihatkan kekuatan teori ukuran
dalam menyerang persoalan-persoalan umum dalam analisis riil dan analisis
fungsional. Namun, semuanya memiliki penerapan di luar teori ukuran, dan
semuanya telah memengaruhi perkembangannya. Inilah bagian-bagian teori
ukuran yang mungkin digunakan oleh setiap analis.
 
Setiap topik idealnya sudah pernah dijumpai oleh mahasiswa sarjana, tetapi
panjang jilid ini memperjelas bahwa mata kuliah sarjana biasa tidak mungkin
mencakup banyak di antaranya. Jilid ini lebih ditujukan pada tingkat
pascasarjana, dan saya berharap isinya memadai untuk menunjang semua mata
kuliah teori ukuran kecuali yang paling ambisius, walaupun barangkali
terlalu padat untuk digunakan langsung sebagai buku ajar, kecuali jika
bagian-bagian yang akan dibahas dipilih dengan cermat. Jika Anda
menggunakannya untuk mempelajari teori ukuran secara mandiri, saya sangat
menganjurkan pendekatan eklektik: carilah pokok bahasan dan teorema tertentu
yang terasa mengejutkan atau berguna, lalu telusuri mundur dari sana. Tujuan
saya yang lain, tentu saja, adalah menyajikan gagasan-gagasan utama teori
ukuran pada tingkat ini secara lebih lengkap daripada yang mudah
ditemukan dalam satu jilid di tempat lain. Saya tidak dapat mengklaim bahwa
uraian ini `definitif', tetapi saya sungguh merasa bahwa saya mencakup
cukup banyak wilayah dengan cara yang memberikan landasan kukuh untuk
studi lebih lanjut. Seperti dalam Jilid 1, saya biasanya tidak menghindar
dari hasil-hasil `terbaik', seperti syarat Lindeberg untuk Teorema Limit
Pusat (\S274), atau teori hasil kali ruang ukur sembarang (\S251).
Jika saya mengajarkan materi ini kepada mahasiswa dalam program doktoral,
saya lebih suka menerima keterbatasan dalam keluasan mata kuliah daripada
membiarkan mereka tidak mengetahui apa yang dapat dilakukan dalam bidang-
bidang yang dibahas.
 
Topik-topik tersebut berinteraksi dengan cara yang rumit -- salah satu
tujuan buku ini ialah menampilkan hubungan-hubungan di antara mereka.
Tidak ada urutan linear kanonik yang harus diikuti untuk mempelajarinya.
Saya pun tidak menganggap bagan organisasi sangat membantu, terutama
karena untuk suatu bab di bagian selanjutnya mungkin hanya dua atau tiga
paragraf dari sebuah bagian yang diperlukan. Setidaknya saya berusaha
menata materi setiap bagian dalam urutan yang membuat bagian awalnya
berguna secara mandiri. Namun, urutan mempelajari bab-bab tersebut sebagian
besar terserah kepada Anda, mungkin setelah melihat sekilas pendahuluan
masing-masing bab. Saya telah berusaha sebaik mungkin untuk mempertahankan
tingkat penyajian yang kurang lebih sama di seluruh jilid, dengan kadang-
kadang membiarkan tingkat kesulitannya meningkat sepanjang suatu bab atau
bagian, tetapi selalu berusaha kembali ke sesuatu yang seharusnya dapat
ditangani oleh siapa pun yang telah menguasai Jilid 1. (Walaupun barangkali
teorema-teorema utama Bab 26 menuntut kerja lebih keras daripada hasil-
hasil pokok di tempat lain, dan \S286 hanya diperuntukkan bagi mereka yang
paling bertekad.)
 
Saya mengatakan ada tujuh topik, sedangkan Anda akan melihat delapan bab
di hadapan Anda. Hal ini karena Bab 21 agak berbeda dari bab-bab lainnya.
Bab ini merupakan teori ukuran yang paling murni, dan ditempatkan di sini
hanya karena pada bagian-bagian selanjutnya dalam jilid ini terdapat tempat
yang (menurut pandangan saya) teorema-teoremanya hanya masuk akal jika
dipahami dalam terang beberapa konsep abstrak yang tidak terlalu sulit,
tetapi juga tidak jelas dengan sendirinya. Meskipun demikian, wajar untuk
mengatakan bahwa gagasan-gagasan terpenting dalam jilid ini sebenarnya
tidak bergantung pada materi Bab 21.
 
Seperti biasa, menentukan seberapa banyak pengetahuan awal yang harus
diandaikan dalam jilid ini merupakan sebuah teka-teki. Tentu saja saya
menggunakan hasil-hasil Jilid 1 dari risalah ini setiap kali hasil-hasil
tersebut tampak relevan. Namun, saya tidak meragukan bahwa akan ada pembaca
yang telah mempelajari teori elementer dari sumber lain. Asalkan Anda dapat
mengonstruksi ukuran Lebesgue dari prinsip-prinsip pertama dan membuktikan
teorema-teorema konvergensi dasar untuk integral pada ruang ukur sembarang,
Anda semestinya dapat mulai mempelajari jilid ini. Barangkali akan membantu
jika Anda memiliki versi Jilid 1 yang hanya memuat hasil-hasil, karena versi
itu mencakup definisi-definisi terpenting maupun pernyataan teorema-
teoremanya.
 
Ada pula pertanyaan mengenai seberapa banyak materi dari luar teori ukuran
yang diperlukan. Bab 21 menuntut sejumlah teori himpunan yang tidak sepele
(diberikan dalam \S2A1), tetapi gagasan-gagasan yang lebih lanjut hanya
diperlukan untuk contoh-contoh tandingan dalam \S216 dan dapat dilewati
pada awalnya. Persoalannya menjadi lebih pelik dalam Bab 24. Di sini kita
memerlukan beragam hasil dari analisis fungsional, beberapa di antaranya
bergantung pada gagasan-gagasan yang tidak sepele dalam topologi umum.
Untuk memahami materi ini sepenuhnya, tidak ada pengganti bagi mata kuliah
tentang ruang bernorma hingga dan termasuk kajian kekompakan lemah. Namun,
saya
tidak suka mensyaratkan persiapan semacam itu, karena kemungkinan besar
persiapan tersebut sekaligus terlalu banyak dan terlalu sedikit. Terlalu
banyak, karena pada tahap ini saya nyaris tidak menyebut operator linear;
terlalu sedikit, karena saya memang meminta sebagian teori ruang yang tidak
konveks lokal, yang sering ditiadakan dalam mata kuliah pertama mengenai
analisis fungsional. Oleh karena itu, dengan risiko memboroskan kertas, saya
telah menuliskan uraian ringkas mengenai fakta-fakta esensial
(\S\S2A3-2A5).  %\S2A3 \S2A4 \S2A5 
 
\bigskip
     
\noindent{\bf Catatan untuk cetakan kedua, April 2003}
     
\medskip
     
Untuk cetakan kedua jilid ini, saya telah membuat dua koreksi substansial
atas pembuktian-pembuktian yang tidak memadai serta sejumlah besar
perubahan kecil; saya sangat berterima kasih kepada T.D.Austin atas
pembacaannya yang cermat terhadap cetakan pertama. Selain itu, saya telah
menambahkan selusin latihan dan beberapa hasil sederhana yang ternyata
relevan bagi materi dalam jilid-jilid selanjutnya dan secara alami cocok
ditempatkan di sini.
     
Proses revisi berkala atas karya ini telah mendorong saya membuat beberapa
inovasi notasi yang tidak dijelaskan secara eksplisit dalam cetakan-cetakan
awal Jilid 1. Saya percaya sebagian besar pembaca akan segera memahaminya.
Namun, jika Anda menemukan rujukan silang yang membingungkan dan tidak
dapat mencocokkannya dengan apa pun dalam versi Jilid 1 yang Anda gunakan,
barangkali ada baiknya memeriksa halaman-halaman erata di {\tt
http://www.essex.ac.uk/maths/people/fremlin/mterr.htm}.
     
\bigskip

\noindent{\bf Catatan untuk edisi sampul keras, Januari 2010}

\medskip

Untuk edisi baru (`Lulu') jilid ini, saya telah menyingkirkan sejumlah
kesalahan lain; tidak diragukan lagi masih banyak yang tersisa. Terdapat
banyak latihan baru, beberapa teorema baru
%214P, 222L Dini derivates, half of \S234, 244O, 251L, \S266, 272Q, 272W
dan beberapa penataan ulang materi yang bersesuaian.
%\S234
Hasil-hasil baru tersebut sebagian besar merupakan tambahan yang hanya
sedikit memengaruhi struktur karya ini, tetapi terdapat sebuah bagian baru
yang singkat (\S266) mengenai ketaksamaan Brunn-Minkowski.

\bigskip

\noindent{\bf Catatan untuk cetakan kedua edisi sampul keras, April 2016}

\medskip

Seperti biasa, ada sekumpulan kesalahan kecil yang harus diperbaiki serta
berbagai perubahan dan tambahan kecil. Namun, alasan utama saya menerbitkan
versi cetak baru adalah suatu cacat besar dalam pembuktian Teorema Carleson;
di sana, langkah gegabah untuk menyederhanakan argumen {\smc Lacey \& Thiele
00} didasarkan pada kesalahan tingkat sarjana\footnote{Saya sangat
berterima kasih kepada A.Derighetti karena telah menarik perhatian saya
pada hal ini}.
Walaupun kekeliruan tersebut cukup kentara, penyelesaiannya tampaknya
memerlukan penyesuaian pada suatu definisi, dan bukanlah tuntutan yang
wajar bagi suatu seminar pascasarjana, yakni khalayak sasaran materi ini.
Lebih jauh lagi, pembuktian tersebut dimaksudkan sebagai ciri
pembeda bukan hanya jilid ini, melainkan juga risalah ini secara
keseluruhan. Maka, dengan permintaan maaf kepada siapa pun yang terpaksa
mundur setelah berhadapan dengan versi asli, saya menyajikan sebuah revisi
yang saya harap pada dasarnya sahih.

\frnewpage 
 
